\clearpage
\begingroup
\let\clearpage\relax
\let\cleardoublepage\relax
\vspace*{\fill}
\chapter{操作系统概述与系统调用}
\begin{center}
操作系统是「管理硬件、服务程序」的系统软件，同时是硬件与应用之间的抽象层与资源仲裁者。本章建立「内核结构--双模式--系统调用--中断」四块基石。
\end{center}
\vspace*{\fill}
\endgroup
\clearpage

\section{操作系统做什么}

\begin{itemize}
    \item \textbf{资源管理者}：分配 CPU、内存、I/O 设备，解决多程序竞争。
    \item \textbf{抽象提供者}：把磁盘抽象成文件、把 CPU 抽象成进程、把内存抽象成地址空间。
    \item \textbf{控制程序}：保护与隔离，防止程序越权。
\end{itemize}

\subsection{内核结构}
\begin{table}[H]
\centering
\caption{内核结构对比}
\begin{tabulary}{\textwidth}{CCC}
\toprule
结构 & 特点 & 代表 \\
\midrule
单体内核 (Monolithic) & 全部服务在内核态，性能高、耦合紧 & Linux \\
微内核 (Microkernel) & 仅核心机制在内核，服务在用户态 & Minix, QNX \\
混合内核 (Hybrid) & 折中 & Windows, macOS \\
\bottomrule
\end{tabulary}
\end{table}

\section{双模式与特权指令}

CPU 提供至少两种模式：
\begin{itemize}
    \item \textbf{用户态 (user mode)}：应用运行，受限。
    \item \textbf{内核态 (kernel mode)}：OS 运行，可执行特权指令。
\end{itemize}
特权指令（如修改页表基址、开关中断、I/O 指令）只能在核心态执行。模式由状态寄存器中的一位控制。

\begin{equation}
    \text{用户态} \xrightarrow{\text{系统调用 / 中断 / 异常}} \text{内核态} \xrightarrow{\text{返回}} \text{用户态}
\end{equation}

\section{系统调用}

系统调用是用户程序请求内核服务的受控入口，通过 \textbf{trap} 指令陷入内核。

\begin{figure}[H]
\centering
\begin{tikzpicture}[font=\small, node distance=1.6cm]
    \node[draw, rounded corners] (app) {用户程序};
    \node[draw, rounded corners, below=of app] (lib) {库函数 (libc)};
    \node[draw, rounded corners, below=of lib] (trap) {\texttt{trap} 陷入};
    \node[draw, rounded corners, below=of trap] (kernel) {内核服务分发};
    \draw[->] (app) -- (lib);
    \draw[->] (lib) -- (trap);
    \draw[->] (trap) -- (kernel);
\end{tikzpicture}
\caption{系统调用的陷入路径}
\end{figure}

常见系统调用类别：进程控制（\texttt{fork}, \texttt{exec}, \texttt{exit}）、文件（\texttt{open}, \texttt{read}, \texttt{write}）、设备（\texttt{ioctl}）、信息（\texttt{getpid}）、通信（\texttt{pipe}, \texttt{socket}）。

\section{中断、异常与陷阱}

\begin{table}[H]
\centering
\caption{三种进入内核的事件}
\begin{tabulary}{\textwidth}{CCC}
\toprule
类型 & 来源 & 例子 \\
\midrule
中断 (Interrupt) & 异步、外部设备 & 时钟、网卡、键盘 \\
异常 (Exception) & 同步、当前指令 & 除零、缺页、非法指令 \\
陷阱 (Trap) & 同步、主动请求 & 系统调用 \\
\bottomrule
\end{tabulary}
\end{table}

\noindent 处理流程：保存现场（PC、寄存器）→ 根据中断向量跳转处理程序 → 恢复现场并返回。时钟中断是抢占式调度的基础。

\begin{equation}
    T_{\text{context}} = T_{\text{save}} + T_{\text{dispatch}} + T_{\text{restore}}
\end{equation}
